\documentclass[10pt,twocolumn]{article}
\usepackage[margin=0.8in]{geometry}
\usepackage{amsmath,amssymb,booktabs,xcolor}
\usepackage[hidelinks]{hyperref}
\usepackage{caption}
\captionsetup{font=small,labelfont=bf}
% Generated by generate_tables.py -- do not edit by hand.
\newcommand{\HOneRows}{ibm\_marrakesh & 0 0 0 0 0 0 1 3 & 0.500 & 0.71\\
ibm\_fez & 0 1 0 0 0 0 0 0 & 0.125 & 0.83\\
ibm\_kingston & 1 0 0 0 0 0 0 0 & 0.125 & 0.71\\}
\newcommand{\HiddenRows}{ibm\_marrakesh & 6 & 11-21 13:30 & q85 & 16.5 & 2.7\\
ibm\_marrakesh & 7 & 11-21 13:46 & q16 & 24.6 & 0.9\\
ibm\_marrakesh & 7 & 11-21 13:46 & q111 & 31.7 & 2.5\\
ibm\_marrakesh & 7 & 11-21 13:46 & q126 & 15.6 & 0.8\\
ibm\_fez & 1 & 11-21 12:08 & q93 & 17.8 & 2.1\\
ibm\_kingston & 0 & 11-20 22:28 & q107 & 22.2 & 1.0\\}
\newcommand{\HTwoEver}{6}
\newcommand{\HTwoRec}{6}
\newcommand{\HThreeFF}{87.4}
\newcommand{\HThreeOF}{0.63}
\newcommand{\HThreeTab}{3169 & 20 \\ fault & 11 & 76}
\newcommand{\SensFF}{88.0}
\newcommand{\SensOF}{0.51}
\newcommand{\NoiseWithin}{0.59}
\newcommand{\NoiseBetween}{1.67}
\newcommand{\ROSevenRows}{ibm\_marrakesh & q93 & disagree & 37.5 & 3.8 & 8.1 & 7.5 & 7.9 & -0.59 [-1.19, +0.01] & +0.39 [-0.33, +1.11]\\
ibm\_marrakesh & q31 & disagree & 37.8 & 4.6 & 24.3 & 24.0 & 24.5 & -0.24 [-1.36, +0.89] & +0.42 [-1.27, +2.11]\\
ibm\_marrakesh & q111 & disagree & 31.7 & 2.5 & 2.4 & 2.0 & 2.1 & -0.37 [-0.56, -0.19] & +0.05 [-0.10, +0.21]\\
ibm\_marrakesh & q130 & agree & 48.9 & 49.0 & 50.0 & 48.8 & 50.2 & -1.16 [-1.99, -0.33] & +1.36 [+0.24, +2.48]\\
ibm\_marrakesh & q48 & agree & 44.5 & 45.8 & 47.6 & 47.9 & 47.4 & +0.30 [-0.45, +1.05] & -0.55 [-1.59, +0.48]\\
ibm\_marrakesh & q14 & control & 0.2 & 0.3 & 0.3 & 0.2 & 0.2 & -0.09 [-0.24, +0.07] & +0.02 [-0.11, +0.14]\\
ibm\_marrakesh & q10 & control & 0.3 & 0.7 & 0.4 & 0.3 & 0.4 & -0.05 [-0.16, +0.06] & +0.07 [-0.08, +0.21]\\
\addlinespace
ibm\_fez & q11 & disagree & 10.1 & 2.7 & 8.3 & 8.8 & 9.7 & +0.43 [-0.20, +1.06] & +0.99 [+0.01, +1.97]\\
ibm\_fez & q71 & disagree & 5.1 & 0.3 & 0.4 & 0.4 & 0.4 & -0.00 [-0.18, +0.17] & -0.04 [-0.15, +0.07]\\
ibm\_fez & q33 & disagree & 7.4 & 3.6 & 1.7 & 1.6 & 1.9 & -0.03 [-0.25, +0.19] & +0.24 [+0.02, +0.45]\\
ibm\_fez & q148 & agree & 19.4 & 20.7 & 19.7 & 19.6 & 19.2 & -0.07 [-0.56, +0.43] & -0.43 [-1.50, +0.64]\\
ibm\_fez & q134 & control & 0.3 & 0.6 & 0.4 & 0.4 & 0.5 & -0.01 [-0.13, +0.11] & +0.06 [-0.11, +0.22]\\
ibm\_fez & q23 & control & 0.3 & 1.1 & 5.2 & 4.9 & 4.6 & -0.33 [-0.59, -0.07] & -0.22 [-0.61, +0.17]\\
\addlinespace
ibm\_kingston & q121 & disagree & 24.3 & 17.1 & 19.5 & 18.9 & 20.5 & -0.64 [-1.54, +0.25] & +1.61 [+0.59, +2.63]\\
ibm\_kingston & q2 & disagree & 6.9 & 1.3 & 0.9 & 1.1 & 1.0 & +0.19 [-0.00, +0.38] & -0.10 [-0.25, +0.05]\\
ibm\_kingston & q21 & disagree & 6.0 & 0.7 & 1.6 & 1.7 & 1.8 & +0.05 [-0.20, +0.29] & +0.08 [-0.16, +0.33]\\
ibm\_kingston & q146 & agree & 48.7 & 49.0 & 50.7 & 49.9 & 50.6 & -0.81 [-2.27, +0.65] & +0.70 [-0.12, +1.52]\\
ibm\_kingston & q148 & control & 0.3 & 0.3 & 0.4 & 0.4 & 0.4 & -0.02 [-0.15, +0.11] & +0.01 [-0.10, +0.13]\\
ibm\_kingston & q93 & control & 0.3 & 0.7 & 0.4 & 0.4 & 0.4 & +0.02 [-0.19, +0.23] & -0.02 [-0.19, +0.16]\\}
\newcommand{\MaxD}{1.61}
\newcommand{\ROSevenJobs}{\texttt{dc0qjqdhdvas73bjtgr0}, \texttt{dc0qk42nj2vs73dee020}, \texttt{dc0qkdeq4uuc73a30sf0}}
\newcommand{\CadRows}{ibm\_marrakesh & 155 0 0 0 0 0 0 0\\
ibm\_fez & 153 0 0 0 0 0 0 153\\
ibm\_kingston & 0 0 0 0 0 0 0 0\\}
\newcommand{\AlignRows}{0 & 11-20 22:28 & 0.881 & 0.20\\
1 & 11-21 12:08 & 0.882 & 0.16\\
2 & 11-21 12:30 & 0.867 & 0.16\\
3 & 11-21 12:45 & 0.899 & 0.15\\
4 & 11-21 13:00 & 0.902 & 0.14\\
5 & 11-21 13:15 & 0.935 & 0.13\\
6 & 11-21 13:30 & 0.953 & 0.12\\
7 & 11-21 13:48 & 0.952 & 0.12\\}
\newcommand{\OldRho}{0.813}
\newcommand{\OldGap}{0.21}
\newcommand{\TotalSeconds}{193}
\newcommand{\NJobs}{27}


\title{\bfseries Stale snapshots, not hidden physics: readout faults on three IBM Heron r2 processors
are rare, persist for tens of minutes, and do not depend on measurement context}
\author{Devin Phillip Davis\\ \small Agile Defense Systems LLC}
\date{21 November 2026 \quad (TESSERACT-RO-06 and RO-07 results, version 1.0)}

\begin{document}
\maketitle

\begin{abstract}
\noindent
Earlier experiments in this series (MAP-03 to MAP-05) found qubits on \texttt{ibm\_marrakesh} whose
in-job readout error reached 34--54\% while IBM's calibration listed about 1\%, and qubit-selection
rules that reversed their ranking within 90 minutes. We ran two preregistered follow-ups
(\NJobs{} jobs, \TotalSeconds{}\,s QPU). In RO-06, an all-qubit readout probe ran 8 times on
\texttt{ibm\_marrakesh}, \texttt{ibm\_fez} and \texttt{ibm\_kingston}. ``Hidden faults'' (probe
$>15\%$, IBM $<3\%$) were rare: 0.13--0.5 per device per round, so H1 was not supported. All
\HTwoEver{} qubits that were ever hidden faults recovered in another round (H2 supported). A faulty
qubit was still faulty 15 minutes later in \HThreeFF\% of cases, against \HThreeOF\% for a healthy
one. RO-07 tested whether readout error depends on context (isolated against simultaneous readout of
all qubits, other qubits in $|0\rangle$ or $|1\rangle$). It does not: every context effect was below
\MaxD\,pp, and both preregistered context hypotheses failed. A further analysis, not preregistered,
points to the explanation. IBM's reported readout errors did not change once in 1.75 hours of
15-minute sampling on any device; \texttt{ibm\_kingston}'s did not change for more than 15 hours.
When \texttt{ibm\_fez}'s values did update, they matched our most recent probes better than the old
values did (Spearman $\rho=0.95$ against $0.81$). The gap between IBM's numbers and in-job readout is
explained by infrequent calibration combined with qubits that change on an hour scale, not by a
difference in how readout is measured.
\end{abstract}

\section{Background}
MAP-03~\cite{map03} reported q46 on \texttt{ibm\_marrakesh} at 54\% readout error against IBM's 1.1\%,
which reversed a 10\,pp layout effect. MAP-05~\cite{map05} saw the same qubit fall to 0.6\% within 90
minutes, and the probe-versus-calibration layout contrast flipped sign on three devices. Two
explanations were open: IBM's calibration and our probe measure readout differently (context), or
IBM's calibration is simply out of date (staleness). RO-06~\cite{pre6} and RO-07~\cite{pre7}, both
preregistered on Zenodo before any job, test the first and describe the time course.

\section{Methods}
\paragraph{RO-06.} Each of 8 rounds submitted one probe job per device (all-$|0\rangle$ and
all-$|1\rangle$ on every qubit, 3 repeats, 2{,}048 shots) and recorded IBM's calibration at
submission. $e_{\mathrm{probe}}$ is the mean of $p(1|0)$ and $p(0|1)$. A hidden fault is
$e_{\mathrm{probe}}>0.15$ with IBM's value $<0.03$, and recovery is $e_{\mathrm{probe}}<0.05$. Rounds
were 15 minutes apart, with one deviation: the client machine slept after round 0 (20 November,
22:28\,UTC), and rounds 1--7 ran from 12:08\,UTC the next day. The 0$\to$1 pair therefore spans
13.5\,h. We report the preregistered analysis together with a sensitivity check that drops it.

\paragraph{RO-07.} On each device, a selection rule fixed in advance and applied to RO-06's last round
chose up to 3 DISAGREE qubits (largest $e_{\mathrm{probe}}-e_{\mathrm{IBM}}$), up to 2 AGREE qubits
(bad readout that IBM reports) and 2 controls, none of them coupled to each other. One job ran 6
interleaved repeats of six circuits: ISO (only the selected qubits measured), ALL (every qubit
prepared alike and measured) and FLIP (selected qubits in $b$, all others in $1-b$, all measured), each
for $b\in\{0,1\}$. The multiplexing effect is $D_{\mathrm{mux}}=e(\mathrm{ALL})-e(\mathrm{ISO})$ and the
neighbour effect is $D_{\mathrm{nbr}}=e(\mathrm{FLIP})-e(\mathrm{ALL})$, each tested with a paired $t$
over repeats (df $=5$). RO-07 ran at about 14:27\,UTC, roughly 40 minutes after RO-06's last round.

\section{Results}

\begin{table}[t]
\centering\scriptsize\setlength{\tabcolsep}{3pt}
\caption{RO-06 H1: hidden faults per round (rounds 0--7) and H4: median Spearman $\rho$ between
probe and IBM readout errors.}
\begin{tabular}{@{}llrr@{}}
\toprule
Device & Hidden faults by round & Mean & $\rho$\\
\midrule
\HOneRows
\bottomrule
\end{tabular}
\end{table}

\begin{table}[t]
\centering\scriptsize\setlength{\tabcolsep}{3pt}
\caption{Every hidden fault observed in RO-06 (readout error, \%).}
\begin{tabular}{@{}lrllrr@{}}
\toprule
Device & Round & UTC & Qubit & Probe & IBM\\
\midrule
\HiddenRows
\bottomrule
\end{tabular}
\end{table}

\paragraph{H1, prevalence: not supported.} No device had a hidden fault in a majority of rounds. Most
badly reading qubits (for example q130 and q48 on \texttt{ibm\_marrakesh} at 45--50\%) are reported
correctly by IBM.

\paragraph{H2, intermittency: supported.} All \HTwoRec{} of \HTwoEver{} qubits that were hidden faults
in some round had recovered ($<5\%$) in another.

\paragraph{H3, persistence.} Over consecutive rounds, P(fault at $r+1$ | fault at $r$) $=$
\HThreeFF\% and P(fault at $r+1$ | healthy at $r$) $=$ \HThreeOF\% (Fisher exact
$p<10^{-100}$). Dropping the 13.5\,h pair gives \SensFF\% and \SensOF\%. Faults typically last at
least tens of minutes, and they appear and disappear between rounds. The median spread between repeats
within a round was \NoiseWithin\,pp, against a median change of \NoiseBetween\,pp between rounds for
flagged qubits, so these changes are not sampling noise.

\begin{table*}[t]
\centering\scriptsize\setlength{\tabcolsep}{3pt}
\caption{RO-07. Readout error (\%) at selection (RO-06 last round: probe and IBM) and in each RO-07
context; $D_{\mathrm{mux}}$ and $D_{\mathrm{nbr}}$ in pp with 95\% CI over 6 repeats. Jobs:
\ROSevenJobs.}
\begin{tabular}{@{}lllrrrrrll@{}}
\toprule
Device & Qubit & Group & Probe & IBM & ISO & ALL & FLIP & $D_{\mathrm{mux}}$ & $D_{\mathrm{nbr}}$\\
\midrule
\ROSevenRows
\bottomrule
\end{tabular}
\end{table*}

\paragraph{RO-07, context: R1 and R2 not supported; R3 (controls) satisfied.} No qubit changed by more
than \MaxD\,pp between contexts, far below the 5\,pp criterion. A qubit at 24\% reads at 24\% whether
measured alone or with 155 others, and whatever the others' states. A few effects are statistically
non-zero at the 1--2\,pp level (q130 on \texttt{ibm\_marrakesh}, q121 on \texttt{ibm\_kingston}, where
$D_{\mathrm{nbr}}$ is positive), which is small multiplexed-readout crosstalk of the size reported in
the literature. Context does not explain the 10--30\,pp gaps. Five of the nine DISAGREE qubits had
also recovered (below 5\%) in the 40 minutes between selection and RO-07 (for example q111 on
\texttt{ibm\_marrakesh}, 31.7\% to 2.0\%), which is another instance of H2.

\begin{table}[t]
\centering\scriptsize\setlength{\tabcolsep}{3pt}
\caption{Exploratory. Qubits (of 156) whose IBM-reported readout error changed between consecutive
snapshots: RO-06 rounds 0$\to$1 (13.5\,h), 1$\to$2, \dots, 6$\to$7 (15\,min each), then 7$\to$RO-07
(40\,min).}
\begin{tabular}{@{}l@{\hspace{4pt}}l@{}}
\toprule
Device & Qubits changed per interval\\
\midrule
\CadRows
\bottomrule
\end{tabular}
\end{table}

\begin{table}[t]
\centering\scriptsize\setlength{\tabcolsep}{3pt}
\caption{Exploratory. \texttt{ibm\_fez}'s new IBM values (updated between 13:48 and 14:26\,UTC)
against the probe from each RO-06 round. Old values against the round-7 probe: $\rho=\OldRho$,
median $|$gap$|$ \OldGap\,pp.}
\begin{tabular}{@{}rlrr@{}}
\toprule
Round & UTC & $\rho$ & Median $|$gap$|$ (pp)\\
\midrule
\AlignRows
\bottomrule
\end{tabular}
\end{table}

\paragraph{Exploratory: calibration cadence.} IBM's readout errors change in discrete updates.
\texttt{ibm\_marrakesh} and \texttt{ibm\_fez} updated overnight. During the 1.75\,h of 15-minute
sampling no device updated, and \texttt{ibm\_kingston} did not update at all between 22:28 on 20
November and 14:26 on 21 November. When \texttt{ibm\_fez} updated, the new values agreed best with
our latest probes (rounds 6--7, $\rho\approx0.95$) and better than the values they replaced
($\rho=\OldRho$). That is what a correct but infrequent measurement would produce.

\section{Discussion}
Taken together, the results point to a simple account. IBM's per-qubit readout errors are accurate
when measured, but they are refreshed only every several to 15 or more hours. Individual qubits move
between about 1\% and 30--50\% readout error on a time scale of tens of minutes to hours. A layout
chosen from the calibration therefore inherits whatever changed since the last update, which accounts
for MAP-03's reversal and MAP-05's sign flip. Context-dependent readout, the alternative mechanism,
is ruled out at the 5\,pp level. The practical rule from MAP-04 and MAP-05 stands: probe all qubits
immediately before a job, and exclude any qubit that is grossly faulty at that moment.

\paragraph{Limitations.} One day; three devices; 8 rounds. The cadence and alignment analyses were not
preregistered. The 15-minute grid cannot resolve faults shorter than about 15 minutes. RO-07 targeted
qubits whose faults had partly recovered by run time, which weakens the context test for the
remaining faulty qubits (q31 and q93 on \texttt{ibm\_marrakesh}, q11 on \texttt{ibm\_fez}, q121 on
\texttt{ibm\_kingston}). We measured only all-0/all-1 and complementary-state contexts.

\section*{Data and code availability}
The record contains all probe counts and calibration snapshots for the 24 RO-06 jobs and 3 RO-07
jobs, the RO-07 locks and QPYs, the per-job usage, the analysis JSONs, and the code that produces
every number here, with a SHA-256 manifest: in \texttt{phase3/}, \texttt{readout\_drift\_06.py},
\texttt{readout\_context\_07.py} and \texttt{readout\_staleness\_explore.py}; and the paper's
\texttt{generate\_tables.py}.

\section*{Acknowledgements}
We acknowledge the use of IBM Quantum services; the views expressed are the author's and do not
reflect the official policy or position of IBM or the IBM Quantum team. Analysis code and this
manuscript were prepared with the assistance of AI coding tools under the author's direction and
review.

\begin{thebibliography}{9}\small
\bibitem{map03} D.~P.~Davis, Zenodo, doi:10.5281/zenodo.24406525 (2026).
\bibitem{map05} D.~P.~Davis, Zenodo, doi:10.5281/zenodo.24407875 (2026).
\bibitem{pre6} D.~P.~Davis, ``Preregistration: TESSERACT-RO-06,'' Zenodo, doi:10.5281/zenodo.24407897 (2026).
\bibitem{pre7} D.~P.~Davis, ``Preregistration: TESSERACT-RO-07,'' Zenodo, doi:10.5281/zenodo.24419327 (2026).
\end{thebibliography}
\end{document}
